\section{Method}
\label{sec:system}

\subsection{Modular Rewrite}

\begin{table}[t]
\centering
\small
\begin{tabularx}{\linewidth}{lX}
\toprule
Module & Responsibility \\
\midrule
\code{envs} & Environment reset, transition, reward, and action parsing. \\
\code{agents} & Model loading, single-request generation, and batched generation. \\
\code{rl\_algos} & PPO and GRPO losses, advantages, and optimizer updates. \\
\code{ragen\_core} & Rollout collection, filtering, batching, and training loops. \\
\code{evaluation} & Episode-level metrics and aggregation. \\
\bottomrule
\end{tabularx}
\caption{Five-module design of \method.}
\label{tab:modules}
\end{table}

The codebase is organized around five modules with explicit contracts. This design makes the six-run matrix a configuration-level controlled comparison. PPO and GRPO share the same environment implementation, agent interface, rollout machinery, logging, evaluation code, and seed. They differ mainly in the RL algorithm module and in whether a critic is present.

\subsection{RAGEN-Compatible Rollout}

The environment builds the system prompt from an environment instruction, grid vocabulary, and action list. Each user turn appends the required output format, including \code{<think>...</think><answer>...</answer>}, and the number of actions remaining. FrozenLake and Sokoban allow a single LLM response to contain multiple atomic actions separated by \code{||}; the base environment executes these actions until termination, truncation, or the action budget is exhausted.

Training and evaluation use the same batched rollout core. For each prompt, multiple independent environment instances are advanced turn by turn, and the agent generates responses for all still-active trajectories in one batch. Completed trajectories are removed from later generation batches. This preserves the per-sequence sampling distribution while improving rollout throughput in native PyTorch.

\subsection{Consumer-Hardware Optimization}

The implementation is designed around a single 8 GB GPU. The key memory concessions are AdamW8bit optimizer states with fp32 embedding states, bf16 model weights, gradient checkpointing, micro-batch size 1 with gradient accumulation 32, and optional reference-policy construction for KL regularization. The most important throughput fix is KV-cache restoration: after HuggingFace gradient checkpointing disables \code{config.use\_cache}, the code restores cache use for rollout generation and explicitly disables cache only inside training forward passes.

These engineering choices are not the main experimental variables; they are the feasibility layer that makes the reproduction possible under consumer hardware. A detailed optimization table is kept in Appendix~\ref{sec:app_impl}.
